\chapter*{Preface}
Tome~4 is about calculating: finding the exact value of a sum, a product or an integral, and, when no exact value
exists, an approximation whose error one can bound. It starts with finite sums, which have a calculus of their own,
and generating functions, which turn sequences into functions. It meets Euler's evaluations of $\sum\frac1{n^2}$ and
of the product for the sine, the integration tricks of the physicists with their justification, and the Gamma
function. Three chapters follow one idea, Fourier analysis, in three settings: the circle, the finite group
$\Z/n\Z$, where it becomes the fast Fourier transform, and the line, where it becomes Poisson summation and decides
the sign of a Gauss sum. Complex analysis then turns real integrals and sums into residues, and leads to the zeta
function and its functional equation. The tome ends with asymptotics, the art of being precise about being
approximate, and with Gauss's arithmetic--geometric mean, which computes $\pi$ with digits doubling at each step.

It assumes Tomes~0 to~3 (references such as ``Tome~0, Theorem~6.4'' point there): mostly series, integrals and
Taylor's formula from Tome~0, arithmetic functions and Gauss sums from Tome~1, characters of finite abelian groups
from Tome~2, and the dominated convergence theorem from Tome~3, Chapter~5.
Exercises are graded \textbf{$\star$} warm-up, \textbf{$\star\star$} for a good afternoon, \textbf{$\star\star\star$}
gems, and every exercise has a hint and a complete solution at the end of its chapter. Each chapter has a
companion Python notebook. Calculation is where a computer is the best of companions: every identity in this book
is checked there to fifty digits, some are first discovered from their decimals and only then proved, and the
notebooks follow the methods into computer science, from the cost of loops and the solution of recurrences to the
fast Fourier transform, big-integer arithmetic and the computation of $\pi$.

The rules of the earlier tomes hold here too. Every statement is proved, or, if quoted only for culture, labelled
as a remark that says where it is proved. Every solution starts with \emph{How to find it} before the
\emph{Write-up}. ``Pause and think'' boxes ask you to guess before you read on. The only fact used without proof is
the one already admitted in Tome~3, the Fubini--Tonelli theorem for general integrands; it is needed rarely, since
Chapter~4 proves Fubini's theorem for continuous functions on rectangles, and each use is labelled
\emph{Admitted (proved in Tome~6)}.

The chapters follow the plan of the series. Chapters~1--3 compute sums, series and products; Chapters~4 and~5
compute integrals; Chapters~6--8 are Fourier analysis; Chapters~9 and~10 are complex analysis and the zeta function;
Chapters~11 and~12 are asymptotics and the arithmetic--geometric mean.
